\documentclass[11pt]{article}
\usepackage[margin=1in]{geometry}
\usepackage{amsmath, amssymb, amsthm}
\usepackage{hyperref}
\usepackage{tocloft}

\title{Homework Problem Sets: Homogeneous Linear Equations}
\author{Curated Practice Problems}
\date{}

\begin{document}

\maketitle
\tableofcontents
\newpage

% ---------------------------------------------------------
\section{Concept Checklist \& Question Mapping}
% ---------------------------------------------------------
To ensure comprehensive coverage of all topics, tricks, and problem archetypes found in the source material, the following concepts have been mapped to the 30 curated problems:

\begin{itemize}
    \item \textbf{C1: General Solution - Distinct Real Roots} (Algebraic factoring, $b^2-4ac > 0$) $\rightarrow$ Q1, Q2, Q3
    \item \textbf{C2: General Solution - Repeated Real Roots} (Perfect square trinomials, multiplying by $t$) $\rightarrow$ Q4, Q5, Q6
    \item \textbf{C3: General Solution - Complex Roots} (Euler's formula, $b^2-4ac < 0$) $\rightarrow$ Q7, Q8, Q9
    \item \textbf{C4: Initial Value Problems (IVP) - Distinct Real Roots} (Solving $2 \times 2$ systems, $e^0=1$ shortcut) $\rightarrow$ Q10, Q11, Q12
    \item \textbf{C5: Initial Value Problems (IVP) - Repeated Real Roots} (Product rule trap for $te^{rt}$) $\rightarrow$ Q13, Q14, Q15
    \item \textbf{C6: Initial Value Problems (IVP) - Complex Roots} (Chain rule and product rule with trig functions) $\rightarrow$ Q16, Q17
    \item \textbf{C7: Boundary Value Problems (BVP)} (Checking for unique, none, or infinite solutions) $\rightarrow$ Q18, Q19, Q20
    \item \textbf{C8: Linear Independence \& The Wronskian} ($2 \times 2$ determinant, $W \neq 0$) $\rightarrow$ Q21, Q22
    \item \textbf{C9: Higher-Order Linear Equations} (Factoring cubics/quartics, $n \times n$ systems) $\rightarrow$ Q23, Q24, Q25
    \item \textbf{C10: Reduction of Order} (Proving the repeated root $t$ multiplier by guessing $y_2 = v(t)y_1$) $\rightarrow$ Q26, Q27
    \item \textbf{C11: Diagnostic Testing ("Find the Flaw")} (Identifying subtle algebraic, calculus, or conceptual errors) $\rightarrow$ Q28, Q29, Q30
\end{itemize}

\newpage

% ---------------------------------------------------------
\section{Problem Sets}
% ---------------------------------------------------------

\subsection*{Part A: General Solutions (Real and Complex Roots)}
\textbf{Directions:} Find the general solution to the following second-order differential equations.
\begin{enumerate}
    \item $y'' - 5y' + 4y = 0$
    \item $y'' + 3y' - 10y = 0$
    \item $2y'' - 7y' + 3y = 0$
    \item $y'' - 8y' + 16y = 0$
    \item $9y'' + 6y' + y = 0$
    \item $y'' + 14y' + 49y = 0$
    \item $y'' + 25y = 0$
    \item $y'' - 4y' + 13y = 0$
    \item $y'' + 2y' + 2y = 0$
\end{enumerate}

\subsection*{Part B: Initial Value Problems (IVPs)}
\textbf{Directions:} Solve the following initial value problems.
\begin{enumerate}
    \setcounter{enumi}{9}
    \item $y'' - y' - 2y = 0; \quad y(0) = 1, \quad y'(0) = 4$
    \item $y'' - 9y = 0; \quad y(0) = 2, \quad y'(0) = -3$
    \item $y'' + 5y' + 6y = 0; \quad y(0) = 0, \quad y'(0) = 2$
    \item $y'' - 4y' + 4y = 0; \quad y(0) = 3, \quad y'(0) = 5$
    \item $y'' + 6y' + 9y = 0; \quad y(0) = 1, \quad y'(0) = -1$
    \item $4y'' - 12y' + 9y = 0; \quad y(0) = 2, \quad y'(0) = 4$
    \item $y'' + 16y = 0; \quad y(0) = 2, \quad y'(0) = -8$
    \item $y'' - 2y' + 5y = 0; \quad y(0) = 1, \quad y'(0) = 3$
\end{enumerate}

\subsection*{Part C: Boundary Value Problems (BVPs)}
\textbf{Directions:} Determine if the following boundary value problems have a unique solution, no solution, or infinitely many solutions. If a solution exists, find it.
\begin{enumerate}
    \setcounter{enumi}{17}
    \item $y'' + 4y = 0; \quad y(0) = 1, \quad y(\pi/4) = 2$
    \item $y'' + 9y = 0; \quad y(0) = 2, \quad y(\pi) = 1$
    \item $y'' + y = 0; \quad y(0) = 3, \quad y(2\pi) = 3$
\end{enumerate}

\subsection*{Part D: Linear Independence and the Wronskian}
\textbf{Directions:} Use the Wronskian or trigonometric identities to determine if the following pairs of functions are linearly independent or dependent on the interval $(0, 1)$.
\begin{enumerate}
    \setcounter{enumi}{20}
    \item $y_1(t) = e^{3t}, \quad y_2(t) = e^{-t}$
    \item $y_1(t) = \sin^2(t), \quad y_2(t) = 1 - \cos(2t)$
\end{enumerate}

\subsection*{Part E: Higher-Order Linear Equations}
\textbf{Directions:} Find the general solution or solve the IVP for the following higher-order equations.
\begin{enumerate}
    \setcounter{enumi}{22}
    \item $y''' - 2y'' - 3y' = 0$
    \item $y''' - y'' = 0; \quad y(0) = 1, \quad y'(0) = 2, \quad y''(0) = 3$
    \item $y^{(4)} - 16y = 0$
\end{enumerate}

\subsection*{Part F: Reduction of Order}
\textbf{Directions:} Given one solution $y_1(t)$, use the reduction of order method (guessing $y_2(t) = v(t)y_1(t)$) to find the second linearly independent solution $y_2(t)$.
\begin{enumerate}
    \setcounter{enumi}{25}
    \item $y'' - 6y' + 9y = 0; \quad y_1(t) = e^{3t}$
    \item $y'' + 4y' + 4y = 0; \quad y_1(t) = e^{-2t}$
\end{enumerate}

\subsection*{Part G: Advanced Diagnostic Testing ("Find the Flaw")}
\textbf{Directions:} You are the grader. Each of the following solutions contains a single mathematical or conceptual error. Identify the error, explain why it is wrong, and provide the correct final solution.
\begin{enumerate}
    \setcounter{enumi}{27}
    \item \textbf{Question:} Solve $y'' - 3y' - 4y = 0$. \\
    \textbf{Flawed Solution:} The auxiliary equation is $r^2 - 3r - 4 = 0$. Factoring yields $(r - 4)(r + 1) = 0$. The roots are $r = 4$ and $r = -1$. The general solution is $y(t) = c_1e^{4t} + c_2e^{t}$.
    
    \item \textbf{Question:} Solve the IVP $y'' + 4y' + 4y = 0; \quad y(0) = 1, \quad y'(0) = 2$. \\
    \textbf{Flawed Solution:} Auxiliary equation: $r^2 + 4r + 4 = 0 \implies (r+2)^2 = 0 \implies r = -2$. General solution: $y(t) = c_1e^{-2t} + c_2te^{-2t}$. Applying $y(0) = 1 \implies c_1 = 1$. Differentiating: $y'(t) = -2c_1e^{-2t} - 2c_2te^{-2t}$. Applying $y'(0) = 2 \implies -2(1) - 0 = 2 \implies -2 = 2$. This is a contradiction, so no solution exists.
    
    \item \textbf{Question:} Calculate the Wronskian of $y_1(t) = \cos t$ and $y_2(t) = \sin t$. \\
    \textbf{Flawed Solution:} $y_1'(t) = -\sin t$ and $y_2'(t) = \cos t$. The Wronskian is $W = y_1y_2' + y_1'y_2$. $W = (\cos t)(\cos t) + (-\sin t)(\sin t) = \cos^2 t - \sin^2 t = \cos(2t)$. Since $W$ is not always zero, they are independent.
\end{enumerate}

\newpage

% ---------------------------------------------------------
\section{Detailed Solutions}
% ---------------------------------------------------------

\subsection*{Part A: General Solutions}

\textbf{Solution 1:} 
Step 1: Auxiliary equation is $r^2 - 5r + 4 = 0$. \\
Step 2: Factor: $(r - 1)(r - 4) = 0 \implies r_1 = 1, r_2 = 4$. \\
Step 3: Distinct real roots. $y(t) = c_1e^t + c_2e^{4t}$.

\textbf{Solution 2:} 
Step 1: Auxiliary equation is $r^2 + 3r - 10 = 0$. \\
Step 2: Factor: $(r + 5)(r - 2) = 0 \implies r_1 = -5, r_2 = 2$. \\
Step 3: Distinct real roots. $y(t) = c_1e^{-5t} + c_2e^{2t}$.

\textbf{Solution 3:} 
Step 1: Auxiliary equation is $2r^2 - 7r + 3 = 0$. \\
Step 2: Factor: $(2r - 1)(r - 3) = 0 \implies r_1 = 1/2, r_2 = 3$. \\
Step 3: Distinct real roots. $y(t) = c_1e^{t/2} + c_2e^{3t}$.

\textbf{Solution 4:} 
Step 1: Auxiliary equation is $r^2 - 8r + 16 = 0$. \\
Step 2: Factor: $(r - 4)^2 = 0 \implies r = 4$ (repeated). \\
Step 3: Repeated root requires multiplying by $t$. $y(t) = c_1e^{4t} + c_2te^{4t}$.

\textbf{Solution 5:} 
Step 1: Auxiliary equation is $9r^2 + 6r + 1 = 0$. \\
Step 2: Factor: $(3r + 1)^2 = 0 \implies r = -1/3$ (repeated). \\
Step 3: $y(t) = c_1e^{-t/3} + c_2te^{-t/3}$.

\textbf{Solution 6:} 
Step 1: Auxiliary equation is $r^2 + 14r + 49 = 0$. \\
Step 2: Factor: $(r + 7)^2 = 0 \implies r = -7$ (repeated). \\
Step 3: $y(t) = c_1e^{-7t} + c_2te^{-7t}$.

\textbf{Solution 7:} 
Step 1: Auxiliary equation is $r^2 + 25 = 0$. \\
Step 2: Solve for $r$: $r^2 = -25 \implies r = \pm 5i$. Here $\alpha = 0, \beta = 5$. \\
Step 3: Complex roots. $y(t) = c_1\cos(5t) + c_2\sin(5t)$.

\textbf{Solution 8:} 
Step 1: Auxiliary equation is $r^2 - 4r + 13 = 0$. \\
Step 2: Quadratic formula: $r = \frac{4 \pm \sqrt{16 - 52}}{2} = \frac{4 \pm \sqrt{-36}}{2} = 2 \pm 3i$. Here $\alpha = 2, \beta = 3$. \\
Step 3: $y(t) = e^{2t}(c_1\cos(3t) + c_2\sin(3t))$.

\textbf{Solution 9:} 
Step 1: Auxiliary equation is $r^2 + 2r + 2 = 0$. \\
Step 2: Quadratic formula: $r = \frac{-2 \pm \sqrt{4 - 8}}{2} = -1 \pm i$. Here $\alpha = -1, \beta = 1$. \\
Step 3: $y(t) = e^{-t}(c_1\cos t + c_2\sin t)$.

\subsection*{Part B: Initial Value Problems}

\textbf{Solution 10:} 
Auxiliary: $r^2 - r - 2 = 0 \implies (r - 2)(r + 1) = 0 \implies r = 2, -1$. \\
General Solution: $y(t) = c_1e^{2t} + c_2e^{-t}$. \\
Derivative: $y'(t) = 2c_1e^{2t} - c_2e^{-t}$. \\
Apply IVs: $y(0) = c_1 + c_2 = 1$ and $y'(0) = 2c_1 - c_2 = 4$. \\
Add equations: $3c_1 = 5 \implies c_1 = 5/3$. Substitute back: $5/3 + c_2 = 1 \implies c_2 = -2/3$. \\
Final Solution: $y(t) = \frac{5}{3}e^{2t} - \frac{2}{3}e^{-t}$.

\textbf{Solution 11:} 
Auxiliary: $r^2 - 9 = 0 \implies r = \pm 3$. \\
General Solution: $y(t) = c_1e^{3t} + c_2e^{-3t}$. \\
Derivative: $y'(t) = 3c_1e^{3t} - 3c_2e^{-3t}$. \\
Apply IVs: $c_1 + c_2 = 2$ and $3c_1 - 3c_2 = -3 \implies c_1 - c_2 = -1$. \\
Add equations: $2c_1 = 1 \implies c_1 = 1/2$. Then $c_2 = 3/2$. \\
Final Solution: $y(t) = \frac{1}{2}e^{3t} + \frac{3}{2}e^{-3t}$.

\textbf{Solution 12:} 
Auxiliary: $r^2 + 5r + 6 = 0 \implies (r + 2)(r + 3) = 0 \implies r = -2, -3$. \\
General Solution: $y(t) = c_1e^{-2t} + c_2e^{-3t}$. \\
Derivative: $y'(t) = -2c_1e^{-2t} - 3c_2e^{-3t}$. \\
Apply IVs: $c_1 + c_2 = 0 \implies c_2 = -c_1$. \\
$y'(0) = -2c_1 - 3c_2 = 2 \implies -2c_1 - 3(-c_1) = 2 \implies c_1 = 2$. Thus $c_2 = -2$. \\
Final Solution: $y(t) = 2e^{-2t} - 2e^{-3t}$.

\textbf{Solution 13:} 
Auxiliary: $r^2 - 4r + 4 = 0 \implies (r - 2)^2 = 0 \implies r = 2$. \\
General Solution: $y(t) = c_1e^{2t} + c_2te^{2t}$. \\
Derivative (Product Rule!): $y'(t) = 2c_1e^{2t} + c_2(e^{2t} + 2te^{2t})$. \\
Apply IVs: $y(0) = c_1 = 3$. \\
$y'(0) = 2c_1 + c_2 = 5 \implies 2(3) + c_2 = 5 \implies 6 + c_2 = 5 \implies c_2 = -1$. \\
Final Solution: $y(t) = 3e^{2t} - te^{2t}$.

\textbf{Solution 14:} 
Auxiliary: $r^2 + 6r + 9 = 0 \implies (r + 3)^2 = 0 \implies r = -3$. \\
General Solution: $y(t) = c_1e^{-3t} + c_2te^{-3t}$. \\
Derivative: $y'(t) = -3c_1e^{-3t} + c_2(e^{-3t} - 3te^{-3t})$. \\
Apply IVs: $y(0) = c_1 = 1$. \\
$y'(0) = -3c_1 + c_2 = -1 \implies -3(1) + c_2 = -1 \implies c_2 = 2$. \\
Final Solution: $y(t) = e^{-3t} + 2te^{-3t}$.

\textbf{Solution 15:} 
Auxiliary: $4r^2 - 12r + 9 = 0 \implies (2r - 3)^2 = 0 \implies r = 3/2$. \\
General Solution: $y(t) = c_1e^{1.5t} + c_2te^{1.5t}$. \\
Derivative: $y'(t) = 1.5c_1e^{1.5t} + c_2(e^{1.5t} + 1.5te^{1.5t})$. \\
Apply IVs: $y(0) = c_1 = 2$. \\
$y'(0) = 1.5c_1 + c_2 = 4 \implies 1.5(2) + c_2 = 4 \implies 3 + c_2 = 4 \implies c_2 = 1$. \\
Final Solution: $y(t) = 2e^{1.5t} + te^{1.5t}$.

\textbf{Solution 16:} 
Auxiliary: $r^2 + 16 = 0 \implies r = \pm 4i$. \\
General Solution: $y(t) = c_1\cos(4t) + c_2\sin(4t)$. \\
Derivative: $y'(t) = -4c_1\sin(4t) + 4c_2\cos(4t)$. \\
Apply IVs: $y(0) = c_1 = 2$. \\
$y'(0) = 4c_2 = -8 \implies c_2 = -2$. \\
Final Solution: $y(t) = 2\cos(4t) - 2\sin(4t)$.

\textbf{Solution 17:} 
Auxiliary: $r^2 - 2r + 5 = 0 \implies r = 1 \pm 2i$. \\
General Solution: $y(t) = e^t(c_1\cos(2t) + c_2\sin(2t))$. \\
Derivative: $y'(t) = e^t(c_1\cos(2t) + c_2\sin(2t)) + e^t(-2c_1\sin(2t) + 2c_2\cos(2t))$. \\
Apply IVs: $y(0) = c_1 = 1$. \\
$y'(0) = c_1 + 2c_2 = 3 \implies 1 + 2c_2 = 3 \implies 2c_2 = 2 \implies c_2 = 1$. \\
Final Solution: $y(t) = e^t(\cos(2t) + \sin(2t))$.

\subsection*{Part C: Boundary Value Problems}

\textbf{Solution 18:} 
Auxiliary: $r^2 + 4 = 0 \implies r = \pm 2i$. General Solution: $y(t) = c_1\cos(2t) + c_2\sin(2t)$. \\
Apply $y(0) = 1 \implies c_1\cos(0) + c_2\sin(0) = c_1 = 1$. \\
Apply $y(\pi/4) = 2 \implies 1\cdot\cos(\pi/2) + c_2\sin(\pi/2) = 0 + c_2(1) = 2 \implies c_2 = 2$. \\
Result: \textbf{Unique solution}, $y(t) = \cos(2t) + 2\sin(2t)$.

\textbf{Solution 19:} 
Auxiliary: $r^2 + 9 = 0 \implies r = \pm 3i$. General Solution: $y(t) = c_1\cos(3t) + c_2\sin(3t)$. \\
Apply $y(0) = 2 \implies c_1 = 2$. \\
Apply $y(\pi) = 1 \implies c_1\cos(3\pi) + c_2\sin(3\pi) = -c_1 = 1 \implies c_1 = -1$. \\
Result: Contradiction ($2 \neq -1$). \textbf{No solution exists}.

\textbf{Solution 20:} 
Auxiliary: $r^2 + 1 = 0 \implies r = \pm i$. General Solution: $y(t) = c_1\cos(t) + c_2\sin(t)$. \\
Apply $y(0) = 3 \implies c_1 = 3$. \\
Apply $y(2\pi) = 3 \implies c_1\cos(2\pi) + c_2\sin(2\pi) = c_1 = 3$. \\
Result: Both conditions yield $c_1 = 3$. $c_2$ is completely unrestricted. \textbf{Infinitely many solutions} of the form $y(t) = 3\cos(t) + c_2\sin(t)$.

\subsection*{Part D: Linear Independence and the Wronskian}

\textbf{Solution 21:} 
$W = \det \begin{pmatrix} e^{3t} & e^{-t} \\ 3e^{3t} & -e^{-t} \end{pmatrix} = (e^{3t})(-e^{-t}) - (3e^{3t})(e^{-t}) = -e^{2t} - 3e^{2t} = -4e^{2t}$. \\
Since $W = -4e^{2t} \neq 0$ for all $t$, the functions are \textbf{linearly independent}.

\textbf{Solution 22:} 
Recall the double angle identity: $\cos(2t) = 1 - 2\sin^2(t)$. \\
Rearranging gives $2\sin^2(t) = 1 - \cos(2t)$. \\
Therefore, $y_2(t) = 2y_1(t)$. Since one is a constant multiple of the other, they are \textbf{linearly dependent} (and $W = 0$).

\subsection*{Part E: Higher-Order Linear Equations}

\textbf{Solution 23:} 
Auxiliary: $r^3 - 2r^2 - 3r = 0$. \\
Factor out $r$: $r(r^2 - 2r - 3) = 0 \implies r(r - 3)(r + 1) = 0$. \\
Roots: $r_1 = 0, r_2 = 3, r_3 = -1$. \\
General Solution: $y(t) = c_1e^{0t} + c_2e^{3t} + c_3e^{-t} = c_1 + c_2e^{3t} + c_3e^{-t}$.

\textbf{Solution 24:} 
Auxiliary: $r^3 - r^2 = 0 \implies r^2(r - 1) = 0$. Roots: $r = 0$ (repeated), $r = 1$. \\
General Solution: $y(t) = c_1 + c_2t + c_3e^t$. \\
Derivatives: $y'(t) = c_2 + c_3e^t$, \quad $y''(t) = c_3e^t$. \\
Apply IVs: \\
$y''(0) = c_3 = 3$. \\
$y'(0) = c_2 + c_3 = 2 \implies c_2 + 3 = 2 \implies c_2 = -1$. \\
$y(0) = c_1 + c_3 = 1 \implies c_1 + 3 = 1 \implies c_1 = -2$. \\
Final Solution: $y(t) = -2 - t + 3e^t$.

\textbf{Solution 25:} 
Auxiliary: $r^4 - 16 = 0$. \\
Factor: $(r^2 - 4)(r^2 + 4) = 0 \implies (r - 2)(r + 2)(r^2 + 4) = 0$. \\
Roots: $r = 2, -2, \pm 2i$. \\
General Solution: $y(t) = c_1e^{2t} + c_2e^{-2t} + c_3\cos(2t) + c_4\sin(2t)$.

\subsection*{Part F: Reduction of Order}

\textbf{Solution 26:} 
Guess $y_2(t) = v(t)e^{3t}$. \\
$y_2' = v'e^{3t} + 3ve^{3t}$. \\
$y_2'' = v''e^{3t} + 6v'e^{3t} + 9ve^{3t}$. \\
Substitute into $y'' - 6y' + 9y = 0$: \\
$(v''e^{3t} + 6v'e^{3t} + 9ve^{3t}) - 6(v'e^{3t} + 3ve^{3t}) + 9(ve^{3t}) = 0$. \\
Factor out $e^{3t}$: $e^{3t}(v'' + 6v' + 9v - 6v' - 18v + 9v) = 0$. \\
$e^{3t}(v'') = 0 \implies v'' = 0 \implies v(t) = c_1t + c_2$. \\
Choosing $c_1 = 1, c_2 = 0$ gives $v(t) = t$. Thus, $y_2(t) = te^{3t}$.

\textbf{Solution 27:} 
Guess $y_2(t) = v(t)e^{-2t}$. \\
$y_2' = v'e^{-2t} - 2ve^{-2t}$. \\
$y_2'' = v''e^{-2t} - 4v'e^{-2t} + 4ve^{-2t}$. \\
Substitute into $y'' + 4y' + 4y = 0$: \\
$(v''e^{-2t} - 4v'e^{-2t} + 4ve^{-2t}) + 4(v'e^{-2t} - 2ve^{-2t}) + 4(ve^{-2t}) = 0$. \\
Factor out $e^{-2t}$: $e^{-2t}(v'' - 4v' + 4v + 4v' - 8v + 4v) = 0$. \\
$e^{-2t}(v'') = 0 \implies v'' = 0 \implies v(t) = t$. Thus, $y_2(t) = te^{-2t}$.

\subsection*{Part G: Advanced Diagnostic Testing ("Find the Flaw")}

\textbf{Solution 28:} 
\textbf{Flaw:} The student correctly found the root $r = -1$ but wrote $e^t$ instead of $e^{-t}$ in the final general solution. \\
\textbf{Correct Solution:} $y(t) = c_1e^{4t} + c_2e^{-t}$.

\textbf{Solution 29:} 
\textbf{Flaw:} The student failed to use the product rule when differentiating the second term $c_2te^{-2t}$. The derivative of $te^{-2t}$ is $e^{-2t} - 2te^{-2t}$, not just $-2te^{-2t}$. \\
\textbf{Correct Solution:} $y'(t) = -2c_1e^{-2t} + c_2(e^{-2t} - 2te^{-2t})$. \\
Applying $y'(0) = 2 \implies -2(1) + c_2(1 - 0) = 2 \implies -2 + c_2 = 2 \implies c_2 = 4$. \\
The correct final solution is $y(t) = e^{-2t} + 4te^{-2t}$.

\textbf{Solution 30:} 
\textbf{Flaw:} The student used the wrong formula for the Wronskian. They added the cross terms ($y_1y_2' + y_1'y_2$) instead of subtracting them ($y_1y_2' - y_1'y_2$). \\
\textbf{Correct Solution:} $W = y_1y_2' - y_1'y_2 = (\cos t)(\cos t) - (-\sin t)(\sin t) = \cos^2 t + \sin^2 t = 1$. Since $W = 1 \neq 0$, the functions are linearly independent.

\newpage

% ---------------------------------------------------------
\section{Concept Checklist \& Question Mapping (Reference)}
% ---------------------------------------------------------
\begin{itemize}
    \item \textbf{C1: General Solution - Distinct Real Roots} $\rightarrow$ Q1, Q2, Q3
    \item \textbf{C2: General Solution - Repeated Real Roots} $\rightarrow$ Q4, Q5, Q6
    \item \textbf{C3: General Solution - Complex Roots} $\rightarrow$ Q7, Q8, Q9
    \item \textbf{C4: Initial Value Problems (IVP) - Distinct Real Roots} $\rightarrow$ Q10, Q11, Q12
    \item \textbf{C5: Initial Value Problems (IVP) - Repeated Real Roots} $\rightarrow$ Q13, Q14, Q15
    \item \textbf{C6: Initial Value Problems (IVP) - Complex Roots} $\rightarrow$ Q16, Q17
    \item \textbf{C7: Boundary Value Problems (BVP)} $\rightarrow$ Q18, Q19, Q20
    \item \textbf{C8: Linear Independence \& The Wronskian} $\rightarrow$ Q21, Q22
    \item \textbf{C9: Higher-Order Linear Equations} $\rightarrow$ Q23, Q24, Q25
    \item \textbf{C10: Reduction of Order} $\rightarrow$ Q26, Q27
    \item \textbf{C11: Diagnostic Testing ("Find the Flaw")} $\rightarrow$ Q28, Q29, Q30
\end{itemize}

\end{document}